% !TeX program = xelatex
\documentclass[a4paper,10pt]{article}

% ============================================================================
% សន្លឹកកិច្ចការសិស្ស ស-ខ្មៅ (Student Worksheet - Grade 8 2.4.6a-b ROOT1)
% បរិបទ៖ គ្មានស្លាយ LCD និងគ្មានម៉ាស៊ីនគិតលេខ — សិស្សបំពេញផ្ទាល់លើសន្លឹកកិច្ចការ ២ ទំព័រ A4
% ============================================================================
\usepackage[a4paper, left=0.85cm, right=0.85cm, top=0.70cm, bottom=0.95cm, footskip=0.48cm]{geometry}
\usepackage{amsmath, amssymb}
\usepackage{unicode-math}
\usepackage{fontspec}
\usepackage{array, colortbl}
\usepackage[table]{xcolor}
\usepackage{enumitem}
\usepackage{pifont}
\usepackage{fancyhdr}
\usepackage{lastpage}
\usepackage{tikz}
\usepackage[most]{tcolorbox}
\usetikzlibrary{calc, arrows.meta, decorations.pathreplacing, positioning, babel}

\makeatletter
\let\@expl@@@mark@update@singlecol@structures@@\relax
\makeatother

\providecommand{\NGSFindFonts}{%
  \IfFileExists{C:/Windows/Fonts/cambria.ttc}{%
    \def\NGSSysPath{C:/Windows/Fonts/}%
    \def\NGSUserPath{C:/Users/KIT/AppData/Local/Microsoft/Windows/Fonts/}%
  }{%
    \IfFileExists{./fonts/KhmerOSsiemreap.ttf}{\def\NGSSysPath{./fonts/}\def\NGSUserPath{./fonts/}}{%
    \IfFileExists{../fonts/KhmerOSsiemreap.ttf}{\def\NGSSysPath{../fonts/}\def\NGSUserPath{../fonts/}}{%
    \IfFileExists{../../fonts/KhmerOSsiemreap.ttf}{\def\NGSSysPath{../../fonts/}\def\NGSUserPath{../../fonts/}}{%
    \IfFileExists{../../../fonts/KhmerOSsiemreap.ttf}{\def\NGSSysPath{../../../fonts/}\def\NGSUserPath{../../../fonts/}}{%
    \IfFileExists{../../../../fonts/KhmerOSsiemreap.ttf}{\def\NGSSysPath{../../../../fonts/}\def\NGSUserPath{../../../../fonts/}}{%
    \IfFileExists{../../../../../fonts/KhmerOSsiemreap.ttf}{\def\NGSSysPath{../../../../../fonts/}\def\NGSUserPath{../../../../../fonts/}}{%
    \IfFileExists{../../../../../../fonts/KhmerOSsiemreap.ttf}{\def\NGSSysPath{../../../../../../fonts/}\def\NGSUserPath{../../../../../../fonts/}}{%
      \def\NGSSysPath{../../../../../../../fonts/}\def\NGSUserPath{../../../../../../../fonts/}%
    }}}}}}}%
  }%
}
\NGSFindFonts

\setmathfont[Path={\NGSSysPath}, FontIndex=1, Scale=1.02]{cambria.ttc}
\defaultfontfeatures{Ligatures=TeX}

\setmainfont[
  Path={\NGSUserPath}, Extension=.ttf, Script=Khmer,
  BoldFont={KhmerOSsiemreap}, BoldFeatures={FakeBold=1.8},
  ItalicFont={KhmerOSsiemreap}, ItalicFeatures={FakeSlant=0.22}
]{KhmerOSsiemreap}

\newfontfamily\muol[Path={\NGSUserPath}, Extension=.ttf, Script=Khmer]{KhmerOSmuollight}
\newfontfamily\TimesEn[Path={\NGSSysPath}, Extension=.ttf, BoldFont={timesbd}, ItalicFont={timesi}]{times}

\XeTeXlinebreaklocale "km"
\XeTeXlinebreakskip = 0pt plus 1pt minus 0.1pt

\definecolor{panelgray}{gray}{0.93}
\definecolor{headgray}{gray}{0.85}

\pagestyle{fancy}
\fancyhf{}
\renewcommand{\headrulewidth}{0pt}
\fancyfoot[L]{\fontsize{7.5pt}{9.5pt}\selectfont អនុវិទ្យាល័យព្រះក្រសាំង \en{|} គណិតវិទ្យាថ្នាក់ទី៨ \en{|} មេរៀនទី២៖ ៤.៦. សញ្ញាណឫសការេ និងឫសគូប (គ្រូ៖ ហុក ប៊ុនថែន)}
\fancyfoot[R]{\fontsize{7.8pt}{9.5pt}\selectfont ទំព័រ \en{\thepage/\pageref{LastPage}}}

\setlength{\parindent}{0pt}
\setlength{\parskip}{0pt}
\setlength{\tabcolsep}{3.8pt}

\newcommand{\en}[1]{{\TimesEn #1}}
\newcommand{\enbf}[1]{{\TimesEn\bfseries #1}}
\newcommand{\unit}[1]{\,\text{\en{#1}}}
\newcommand{\cnum}[1]{\mbox{\fontsize{10pt}{12pt}\selectfont\ding{\numexpr181+#1\relax}}}
\newcommand{\chkbox}{\tikz[baseline=-0.12ex]\draw[black, line width=0.8pt] (0,0) rectangle (0.24cm,0.24cm);}
\newcommand{\blank}[1][1.5cm]{\makebox[#1]{\dotfill}}
\newcommand{\ansbox}[1][0.85cm]{\,\tikz[baseline=-0.5ex]\draw[black, line width=0.65pt, fill=white] (0,0) rectangle (#1, 0.42cm);\,}
\newcommand{\secbar}[2]{%
  \begin{tcolorbox}[enhanced, sharp corners, colback=headgray, colframe=black, boxrule=0.9pt, boxsep=1pt, top=2pt, bottom=2pt, left=5pt, right=5pt, before skip=4pt, after skip=3pt]
    {\muol\fontsize{9.2pt}{13pt}\selectfont #1}\hfill{\fontsize{8pt}{11pt}\selectfont\textbf{#2}}
  \end{tcolorbox}%
}

\newcolumntype{L}[1]{>{\raggedright\arraybackslash}p{#1}}
\newcolumntype{C}[1]{>{\centering\arraybackslash}p{#1}}

\begin{document}
\fontsize{9.6pt}{14.4pt}\selectfont

% ============================================================================
% ទំព័រទី ១៖ E1 + E2 (POE 1 & POE 2) + E3 សំយោគនិយមន័យ និងធនាគារហេតុផល
% ============================================================================
\begin{tcolorbox}[enhanced, sharp corners, colback=white, colframe=black, boxrule=1.0pt, boxsep=2pt, top=3pt, bottom=3pt, left=6pt, right=6pt]
  \begin{center}
    {\muol\fontsize{10.3pt}{15pt}\selectfont អនុវិទ្យាល័យព្រះក្រសាំង \enbf{--} សន្លឹកកិច្ចការសិស្ស \enbf{(}មេរៀនទី២៖ ៤\enbf{.}៦\enbf{.} សញ្ញាណឫសការេ និងឫសគូប\enbf{)}}\\[3pt]
    {\fontsize{8.8pt}{12pt}\selectfont
    ឈ្មោះសិស្ស / ដៃគូ៖ \blank[5.2cm] \hspace{0.25cm}
    ថ្នាក់ទី៨៖ \blank[1.1cm] \hspace{0.25cm}
    ក្រុមទី៖ \blank[1.0cm] \hspace{0.25cm}
    កាលបរិច្ឆេទ៖ ១៩/១១/២០២៦}
  \end{center}
\end{tcolorbox}

% --- E1 + E2 ---
\secbar{ផ្នែកទី១៖ បេសកកម្មទស្សន៍ទាយ \enbf{(E1)} និងការរុករកជាដៃគូ \enbf{(E2 -- POE 1 \& POE 2)}}{រយៈពេល ១៤ នាទី}

\begin{tcolorbox}[enhanced, sharp corners, colback=panelgray, colframe=black, boxrule=0.75pt, boxsep=1.5pt, top=2.5pt, bottom=2.5pt, left=5pt, right=5pt]
  \textbf{បេសកកម្ម \en{E1} (គូសធីក \chkbox\ មុនពិសោធន៍)៖}
  (១) តម្លៃនៃនិមិត្តសញ្ញា $\sqrt{16}$ ស្មើ៖ \chkbox\ $+4$ តែមួយ \quad \chkbox\ $\pm 4$ ទាំងពីរ\\
  (២) ក្នុងចំណោម $x^2 = -9$ និង $\sqrt[3]{-8}$៖ \chkbox\ គ្មានឫសទាំងពីរ \quad \chkbox\ $x^2=-9$ គ្មានឫស តែ $\sqrt[3]{-8} = -2$
\end{tcolorbox}

\vspace{3pt}
\textbf{ក. សកម្មភាពទី ១ (\en{POE 1})៖ តារាងទ្វេទិស សមីការ $x^2 = a$ ទល់នឹង និមិត្តសញ្ញារ៉ាឌីកាល់ $\sqrt{a}$ (ទំ.២២)}

{\renewcommand{\arraystretch}{1.50}
\begin{tabular}{|C{2.40cm}|L{4.85cm}|L{4.75cm}|C{2.80cm}|C{2.80cm}|}
  \hline
  \rowcolor{headgray}
  \textbf{តម្លៃ $a$} & \textbf{ពិនិត្យការេ $(\pm k)^2 = a$} & \textbf{ឫសទាំងអស់នៃសមីការ $x^2 = a$} & \textbf{ឫសវិជ្ជមាន $\sqrt{a}$} & \textbf{ឫសអវិជ្ជមាន $-\sqrt{a}$} \\
  \hline
  $a = 16$ & $(+4)^2 = 16$ និង $(-4)^2 = 16$ & $x = +4$ \;ឬ\; $x = -4$ (២ ឫស) & $\sqrt{16} = +4$ & $-\sqrt{16} = -4$ \\
  \hline
  $a = 25$ & $(+5)^2 = \blank[0.9cm]$ , $(-5)^2 = \blank[0.9cm]$ & $x = \blank[1.0cm]$ ឬ $x = \blank[1.0cm]$ & $\sqrt{25} = \blank[1.0cm]$ & $-\sqrt{25} = \blank[1.0cm]$ \\
  \hline
  $a = 0.09$ & $(+0.3)^2 = 0.09$ , $(-0.3)^2 = 0.09$ & $x = \blank[1.0cm]$ ឬ $x = \blank[1.0cm]$ & $\sqrt{0.09} = \blank[0.9cm]$ & $-\sqrt{0.09} = \blank[0.9cm]$ \\
  \hline
  $a = 0$ & $0^2 = 0$ & $x = \blank[1.0cm]$ (តែ ១ ឫសគត់) & $\sqrt{0} = \blank[1.0cm]$ & គ្មានសញ្ញាដក \\
  \hline
  $a = -9$ & គ្មានចំនួនណាលើកជាការេបានដកទេ! & \chkbox\ មានឫស \quad \chkbox\ គ្មានឫស & \multicolumn{2}{C{5.60cm}|}{មិនកំណត់ក្នុងសំណុំចំនួនសនិទាន/ពិត!} \\
  \hline
\end{tabular}}

\vspace{3pt}
$\bullet$ \textbf{ពន្យល់ (\en{E})៖} សញ្ញា $\sqrt{\phantom{x}}$ ហៅថា \blank[2.3cm] ចំណែកចំនួន $a$ ក្រោមរ៉ាឌីកាល់ហៅថា \blank[2.3cm] ហើយ $\sqrt{a}$ មានតម្លៃ \blank[1.6cm] ជានិច្ច ។

\vspace{4pt}
\textbf{ខ. សកម្មភាពទី ២ (\en{POE 2})៖ ឫសគូប $\sqrt[3]{a}$ និងច្រកសញ្ញាផលគុណក្រោមរ៉ាឌីកាល់ (ទំ.២៣–២៤)}

{\renewcommand{\arraystretch}{1.50}
\begin{tabular}{|C{3.60cm}|L{5.65cm}|C{4.15cm}|L{4.20cm}|}
  \hline
  \rowcolor{headgray}
  \textbf{សមីការ / កន្សោម} & \textbf{ពិនិត្យស្វ័យគុណទី ៣ ឬសញ្ញាផលគុណ} & \textbf{តម្លៃឫសត្រឹមត្រូវ} & \textbf{សង្កេតសញ្ញា} \\
  \hline
  $x^3 = +8$ & $(+2)^3 = (+2)(+2)(+2) = +8$ & $\sqrt[3]{8} = +2$ & វិជ្ជមាន $\to$ ឫសវិជ្ជមាន \\
  \hline
  $x^3 = -8$ & $(-2)^3 = (-2)(-2)(-2) = -8$ & $\sqrt[3]{-8} = \blank[1.3cm]$ & អវិជ្ជមាន $\to$ ឫស \blank[1.5cm] \\
  \hline
  $x^3 = -64$ & $(\blank[1.0cm])^3 = -64$ & $\sqrt[3]{-64} = \blank[1.3cm]$ & មានឫសតែ \blank[0.8cm] គត់ \\
  \hline
  $x^3 = 0.008$ & $(+0.2)^3 = 0.008$ & $\sqrt[3]{0.008} = \blank[1.2cm]$ & វិជ្ជមាន $\to$ ឫស \blank[1.5cm] \\
  \hline
  $\sqrt{(-9)(-16)}$ & គិតផលគុណមុន៖ $(-9)(-16) = +144 > 0$ & $= \sqrt{144} = \blank[1.2cm]$ & $(-)\times(-) = (+)$ មានឫស! \\
  \hline
\end{tabular}}

\vspace{3pt}
$\bullet$ \textbf{ពន្យល់ (\en{E})៖} ហេតុអ្វី $x^2 = -9$ គ្មានឫស តែ $x^3 = -27$ មានឫសស្មើ $-3$? ព្រោះផលគុណចំនួនអវិជ្ជមាន $3$ កត្តាមានសញ្ញា \blank[1.8cm] ។

% --- E3 ---
\secbar{ផ្នែកទី២៖ សំយោគក្បួនទូទៅ \enbf{(E3)} និងធនាគារហេតុផល \enbf{B\&W} \cnum{1}\enbf{--}\cnum{5}}{សៀវភៅពុម្ព ទំ.២២–២៣}

\begin{tabular}{|L{9.20cm}|L{9.20cm}|}
  \hline
  \textbf{១. ក្បួនជាទូទៅ និងសម្គាល់ (ទំ.២២–២៣)៖}\par\vspace{2pt}
  $\bullet$ \textbf{ឫសការេ៖} បើ $a > 0$ សមីការ $x^2 = a$ មានឫសការេពីរផ្ទុយគ្នាគឺ $x = +\sqrt{a}$ \textbf{ឬ} $x = -\sqrt{a}$ ។\par\vspace{2pt}
  $\bullet$ បើ $a = 0$ នោះ $\sqrt{0} = \blank[0.7cm]$ ។ បើ $a < 0$ នោះ \textbf{គ្មានឫសការេទេ} ។\par\vspace{2pt}
  $\bullet$ \textbf{ឫសគូប៖} ឫសគូបនៃ $a$ គឺចំនួន $x$ ដែល $x^3 = a$ កំណត់សរសេរ $x = \ansbox[1.0cm]$ (មានតែមួយគត់គ្រប់សញ្ញានៃ $a$) ។
  &
  \textbf{២. ធនាគារហេតុផល \en{B\&W} (\cnum{1}--\cnum{5})៖}\par\vspace{2pt}
  \cnum{1} សមីការ $x^2 = a$ ($a>0$) $\implies$ មានឫសពីរ $x = +\sqrt{a}$ ឬ $x = -\sqrt{a}$ ។\par\vspace{2pt}
  \cnum{2} និមិត្តសញ្ញា $\sqrt{a}$ ($a\ge 0$) $\implies$ តាងឫសការេវិជ្ជមានតែមួយគត់ ($\sqrt{a}\ge 0$) ។\par\vspace{2pt}
  \cnum{3} ច្រកសញ្ញាឫសការេ ($a<0$) $\implies$ ចំនួនអវិជ្ជមានគ្មានឫសការេទេ ។\par\vspace{2pt}
  \cnum{4} ឫសគូប $\sqrt[3]{a}$ $\implies$ មានតែមួយគត់ ហើយរក្សាសញ្ញាដូច $a$ ជានិច្ច ។\par\vspace{2pt}
  \cnum{5} ផលគុណក្រោមរ៉ាឌីកាល់ $\implies$ គិតសញ្ញាផលគុណក្នុងរ៉ាឌីកង់ជាមុនសិន! \\
  \hline
\end{tabular}

\vspace{5pt}
{\renewcommand{\arraystretch}{1.52}
\begin{tabular}{|C{4.40cm}|C{4.40cm}|C{4.40cm}|C{4.40cm}|}
  \hline
  \rowcolor{panelgray}
  $\sqrt{49} = \blank[1.8cm]$ & $-\sqrt{49} = \blank[1.8cm]$ & $\sqrt[3]{1000} = \blank[1.7cm]$ & $\sqrt[3]{-1000} = \blank[1.6cm]$ \\
  \hline
  $\sqrt{0.25} = \blank[1.7cm]$ & $\sqrt{\dfrac{9}{16}} = \blank[1.7cm]$ & $\sqrt[3]{-125} = \blank[1.6cm]$ & $\sqrt{(-4)(-25)} = \blank[1.3cm]$ \\
  \hline
\end{tabular}}

\newpage

% ============================================================================
% ទំព័រទី ២៖ E4 លំហាត់គំរូ + E5 ប្រតិបត្តិ + វិនិច្ឆ័យ A-D + Exit Check + កិច្ចការផ្ទះ
% ============================================================================
\secbar{ផ្នែកទី៣៖ ពង្រីកចំណេះដឹងតាមជំហានស៊ើបអង្កេត \enbf{(E4 --} លំហាត់គំរូ ទំ\enbf{.}២២\enbf{--}២៣\enbf{)}}{រយៈពេល ១០ នាទី}

\textbf{១. លំហាត់គំរូទី ១ (ដោះស្រាយសមីការ ទំ.២២–២៣)៖}

{\renewcommand{\arraystretch}{1.46}
\begin{tabular}{|L{3.10cm}|L{8.80cm}|C{3.75cm}|C{2.15cm}|}
  \hline
  \rowcolor{headgray}
  \textbf{សមីការ} & \textbf{ការវិភាគសញ្ញា និងស្វ័យគុណច្រាស} & \textbf{ចម្លើយ} & \textbf{ហេតុផល} \\
  \hline
  \textbf{ក.} $x^2 = 16$ & $16 > 0 \implies$ មានឫសការេពីរផ្ទុយគ្នា $+\sqrt{16}$ និង $-\sqrt{16}$ & $x = \blank[0.8cm]$ ឬ $x = \blank[0.8cm]$ & \cnum{1} \\
  \hline
  \textbf{ខ.} $x^2 = -9$ & $-9 < 0$ ហើយគ្រប់ការេ $x^2 \ge 0$ ជានិច្ច & \blank[3.0cm] & \cnum{3} \\
  \hline
  \textbf{គ.} $x^3 = 27$ & $27 = 3^3 \implies x = \sqrt[3]{27}$ (ឫសគូបវិជ្ជមាន) & $x = \blank[2.0cm]$ & \cnum{4} \\
  \hline
  \textbf{ឃ.} $x^3 = -27$ & $-27 = (-3)^3 \implies x = \sqrt[3]{-27}$ (ឫសគូបអវិជ្ជមាន) & $x = \blank[2.0cm]$ & \cnum{4} \\
  \hline
\end{tabular}}

\vspace{5pt}
\textbf{២. លំហាត់គំរូទី ២ (ចំណោទធរណីមាត្រការេ និងគូប ទំ.២៣ — សម្គាល់លក្ខខណ្ឌប្រវែងជ្រុង $x > 0$)៖}

\begin{tabular}{|L{9.20cm}|L{9.20cm}|}
  \hline
  \textbf{១) ការេមានក្រឡាផ្ទៃ $S = 121\unit{dm}^2$ ។ រកប្រវែងជ្រុង $x$៖}\par\vspace{3pt}
  $\bullet$ សមីការក្រឡាផ្ទៃ៖ $x^2 = 121$\par\vspace{3pt}
  $\bullet$ ដោយ $x$ ជាប្រវែងជ្រុង ($x > 0$) យើងយកតែឫសវិជ្ជមាន៖\par\vspace{3pt}
  $x = \sqrt{121} = \sqrt{(\blank[0.8cm])^2} = \blank[1.3cm]\unit{dm}$ \quad (\cnum{2})
  &
  \textbf{២) គូបមានមាឌ $V = 125\unit{cm}^3$ ។ រកប្រវែងជ្រុង $x$៖}\par\vspace{3pt}
  $\bullet$ សមីការមាឌគូប៖ $x^3 = 125$\par\vspace{3pt}
  $\bullet$ បំបែក $125 = (\blank[0.8cm])^3$ រួចបញ្ចេញពីឫសគូប៖\par\vspace{3pt}
  $x = \sqrt[3]{125} = \blank[1.3cm]\unit{cm}$ \quad (\cnum{4}) \\
  \hline
\end{tabular}

\vspace{4pt}
$\bullet$ \textbf{សំណួរឆ្លុះបញ្ចាំង \en{(EIA Looking-Back)}៖} ហេតុអ្វីក្នុងចំណោទរកប្រវែងជ្រុងការេ យើងមិនយកឫសអវិជ្ជមាន $x = -11\unit{dm}$? \blank[3.6cm]

% --- E5 ---
\secbar{ផ្នែកទី៤៖ វាយតម្លៃលទ្ធផលសិក្សា \enbf{(E5 --} ប្រតិបត្តិ ទំ\enbf{.}២៣\enbf{--}២៤ \enbf{+} វិនិច្ឆ័យ \enbf{A--D +} សំបុត្រចេញ\enbf{)}}{រយៈពេល ១០ នាទី}

\textbf{ក. លំហាត់ប្រតិបត្តិ (សៀវភៅពុម្ព ទំ.២៣ និង ទំ.២៤ លើ — ធ្វើជាដៃគូ រួចប្ដូរសន្លឹកផ្ទៀង \en{Pause-Review})៖}

{\renewcommand{\arraystretch}{1.50}
\begin{tabular}{|L{8.90cm}|L{8.90cm}|}
  \hline
  \textbf{១.ក (ទំ.២៣)៖} $x^2 = 196 \implies x = \blank[2.8cm]$ (\cnum{1}) &
  \textbf{១.ខ (ទំ.២៣)៖} ការេ $S = 2.25\unit{m}^2 \implies x = \blank[2.0cm]\unit{m}$ (\cnum{2}) \\
  \hline
  \textbf{២.ក (ទំ.២៤)៖} $\sqrt{(-64)(-81)} = \sqrt{64\times 81} = \blank[2.1cm]$ (\cnum{5}) &
  \textbf{២.ខ (ទំ.២៤)៖} $\sqrt{-1225}$៖ \blank[3.4cm] (\cnum{3}) \\
  \hline
  \multicolumn{2}{|L{18.0cm}|}{\textbf{២.គ (ទំ.២៤)៖} $\sqrt[3]{(-128)\times 108} = \sqrt[3]{(-2^7)\times(2^2\times 3^3)} = \sqrt[3]{(-2^3)^3\times 3^3} = (\blank[1.0cm])\times 3 = \blank[1.8cm]$ \quad (\cnum{4})} \\
  \hline
\end{tabular}}

\vspace{5pt}
\textbf{ខ. គូសវង់ជុំវិញជម្រើសត្រឹមត្រូវតែ ១ គត់ ក្នុងចំណោម \en{A, B, C, D} រួចសរសេរលេខហេតុផល \cnum{1}--\cnum{5} នៃជម្រើសខុស៖}

{\renewcommand{\arraystretch}{1.35}
\begin{tabular}{|C{4.40cm}|C{4.40cm}|C{4.40cm}|C{4.40cm}|}
  \hline
  \cellcolor{panelgray}\textbf{ជម្រើស \en{A}} & \cellcolor{panelgray}\textbf{ជម្រើស \en{B}} & \cellcolor{panelgray}\textbf{ជម្រើស \en{C}} & \cellcolor{panelgray}\textbf{ជម្រើស \en{D}} \\
  \hline
  $\sqrt{(-64)(-81)} = 72$ & $\sqrt{16} = \pm 4$ & $x^2 = 196 \Rightarrow x = 14$ តែមួយ & $\sqrt[3]{-27}$ គ្មានឫស \\[3pt]
  \chkbox\ ត្រូវ \quad \chkbox\ ខុស (\blank[0.7cm]) & \chkbox\ ត្រូវ \quad \chkbox\ ខុស (\blank[0.7cm]) & \chkbox\ ត្រូវ \quad \chkbox\ ខុស (\blank[0.7cm]) & \chkbox\ ត្រូវ \quad \chkbox\ ខុស (\blank[0.7cm]) \\
  \hline
\end{tabular}}

\vspace{5pt}
\begin{tcolorbox}[enhanced, sharp corners, colback=white, colframe=black, boxrule=0.85pt, boxsep=2.5pt, top=3.5pt, bottom=3.5pt, left=5pt, right=5pt]
  {\muol\fontsize{8.8pt}{12.5pt}\selectfont សំបុត្រចេញរហ័ស \enbf{(Exit Check --} ៦០ វិនាទី\enbf{)}៖}\\[2pt]
  (១) $\sqrt{81} = \blank[1.8cm]$ \quad \en{|} \quad
  (២) សមីការ $x^2 = 49 \implies x = \blank[2.2cm]$ \quad \en{|} \quad
  (៣) $\sqrt[3]{-64} = \blank[1.8cm]$
\end{tcolorbox}

\vspace{5pt}
\begin{tcolorbox}[enhanced, sharp corners, colback=panelgray, colframe=black, boxrule=0.85pt, boxsep=2.5pt, top=3.5pt, bottom=3.5pt, left=5pt, right=5pt]
  {\muol\fontsize{8.8pt}{12.5pt}\selectfont កិច្ចការផ្ទះ \enbf{(}សៀវភៅពុម្ព ទំព័រ ២៦ លំហាត់ \enbf{\#12, \#13, \#14} និង \enbf{\#15)}៖}\\[2pt]
  $\bullet$ \textbf{លំហាត់ \#12៖} រកប្រវែងជ្រុងការេដែលមានក្រឡាផ្ទៃ៖ ក) $81\unit{cm}^2$ ; ខ) $0.04\unit{m}^2$ ; គ) $\dfrac{16}{25}\unit{dm}^2$ ។\\
  $\bullet$ \textbf{លំហាត់ \#13៖} ដោះស្រាយសមីការ $x^2 = 36$, $x^2 = 121$, $x^2 = 1.44$, $x^2 = -25$ ។ \quad
  $\bullet$ \textbf{\#14៖} $x^3 = 64$, $x^3 = -125$, $x^3 = 0.008$ ។ \quad
  $\bullet$ \textbf{\#15៖} គូបមានមាឌ $216\unit{cm}^3$ ។ រកប្រវែងជ្រុង ។
\end{tcolorbox}

\end{document}
